% !TEX root = main_part4_11.tex
\documentclass[reprint,amsmath,amssymb,aps,prd]{revtex4-2}

\usepackage{graphicx}
\usepackage{dcolumn}
\usepackage{bm}
\usepackage{hyperref}

\begin{document}

\preprint{APS/123-QED}

\title{Mechanics of Spatial Vibration IV: Emergent SU(3) Hadronic Symmetries and Topological Tensor Knots in 3D Geomechanics}

\author{Kwang Yong Yoo}
\email{khanyong.yoo@gmail.com}
\affiliation{Independent Researcher \\
\textit{AI-Assisted Research and Simulation Development Support: Gemini} \\
(Interactive simulations and supplementary materials are officially archived under DOI: \href{https://doi.org/10.5281/zenodo.21438016}{10.5281/zenodo.21438016})}

\date{\today}

\begin{abstract}
The Standard Model characterizes hadrons via SU(3) abstract algebra, attributing these symmetries to unobservable point-like quarks within a mathematical internal space. While phenomenologically successful, this reliance on abstract mathematical constructs leaves the physical, 3-dimensional geometric origins of these symmetries largely unexplored. To bridge this conceptual gap, building strictly upon the \textbf{`Geometrical Fluctuation of Space'} (Part I) and the \textbf{`Spatial Vibration Tensor'} (Part II) framework, this paper proposes that the non-abelian gauge symmetries of Quantum Chromodynamics (QCD) are not fundamental, but rather \textit{emergent} collective behaviors of highly stable \textbf{`Topological Tensor Knots'} formed by vortices within the 3D spatial tensor fluid.

To overcome the heuristic limitations of phenomenological potentials, we analytically approach Color Confinement via the topological conservation of fluid helicity (Hopf invariant). Driven by \textbf{`Tensor Condensation'} (Part II), extending a spatial knot strictly mandates a linearly increasing geometric tension, directly deriving the QCD string tension ($\sigma$) from the dynamics of the \textbf{`geometrical stress tensor'} ($S_{\mu\nu}$) without invoking ad-hoc parameters. Furthermore, we abandon the direct mathematical isomorphism between the 3D real space $SO(3)$ and the 8-dimensional complex space $SU(3)$. Instead, we propose an emergent framework where the non-abelian $SU(3)$ multiplet structures arise as macroscopic collective excitations of intertwined fundamental 3D topological defects.

Addressing the critical inconsistency between discrete topological variables and continuous differential equations, we reformulate the decay of ``Strangeness'' (Weak Interaction) not as continuous viscous damping, but as discontinuous topological jumps governed by quantum instanton tunneling rates ($\Gamma \propto e^{-S_{\mathrm{inst}}/\hbar}$). The Pauli Exclusion Principle is geometrically re-evaluated not as classical hydrodynamic repulsion, but as an infinite topological deformation energy barrier mandated by the interference of $4\pi$-twisted spinor phase planes. By replacing internal dimensions with the emergent topological dynamics of a 3D spatial vacuum, this study offers a mathematically consistent, deterministic realization of hadronic symmetries under the universal principle of \textbf{`Fractal Scale Symmetry'}.
\end{abstract}

\maketitle

\section{Introduction}

\subsection{The Epistemological Limits of Abstract Algebra in Particle Physics}
The Standard Model of particle physics successfully categorizes hadrons using SU(3) abstract algebra \cite{GellMann1964}, attributing these symmetries to unobservable point-like quarks residing within a mathematical internal space. While phenomenologically successful, this reliance on abstract mathematical constructs leaves the physical, 3-dimensional geometric origins of these symmetries largely unexplained. To bridge this conceptual gap, we extend the geomechanical framework developed in the \textit{Mechanics of Spatial Vibration} trilogy \cite{YooPart1, YooPart2, YooPart3} into the subatomic realm. Although Quantum Chromodynamics (QCD) predicts experimental outcomes with high precision, it relies on phenomenological potentials to enforce ``Color Confinement'' \cite{Wilson1974} and abstract Higgs mechanisms for mass generation \cite{Higgs1964}, without elucidating the underlying geometric mechanisms within the physical spatial vacuum.

\subsection{Objectives of the Geomechanical Framework}
Based strictly on the \textbf{`Geometrical Fluctuation of Space'} hypothesis (Part I) \cite{YooPart1} and the phenomenological \textbf{`Spatial Vibration Tensor'} ($\tilde{V}_{\mu\nu}$) framework (Part II) \cite{YooPart2}, this paper proposes a rigorous geomechanical derivation of hadronic symmetries. We propose that hadrons are not composites of fundamental point particles, but rather highly stable \textbf{`Topological Tensor Knots'} formed when extreme-energy collisions fracture the 3D spatial tensor fluid. By bridging abstract group theory with deterministic spatial geomechanics, we aim to demonstrate that the abstract non-abelian symmetries of particle physics are \textit{emergent} properties of a physical, 3-dimensional tensor vacuum.

\section{Beyond Point Particles: Topological Tensor Knots and Confinement}

QCD introduces a phenomenological linear potential ($V(r) \sim \sigma r$) to enforce Color Confinement and observes Asymptotic Freedom at short distances \cite{Gross1973}. We derive both confinement and asymptotic freedom from the principles of topological fluid dynamics.

\subsection{Fluid Helicity and the Hopf Invariant}
When a particle accelerator injects extreme energy into the vacuum, the shockwave violently twists the spatial tensor fluid into self-sustaining standing waves. We define the velocity field of this spatial fluid as $\mathbf{v}$ and its vorticity as $\boldsymbol{\omega} = \nabla \times \mathbf{v}$. To quantify the indivisibility of these knots, we introduce the Hopf invariant ($\mathcal{H}$) \cite{Moffatt1969}, representing the topological fluid helicity:
\begin{equation}
\mathcal{H} = \int \mathbf{v} \cdot \boldsymbol{\omega} \, d^3x = n \cdot h_0 \quad (n \in \mathbb{Z})
\end{equation}
where $h_0$ is the fundamental quantum of spatial fluid helicity, setting the required $[L^4/T^2]$ dimensional scale. Because the topological winding $n$ must remain an integer, severing the knot inevitably forces the fluid to form new topological boundaries, guaranteeing Color Confinement mathematically.

\subsection{Transverse Geometric Pressure and String Tension ($\sigma$)}
Classical fluid dynamics suggests the cross-sectional area ($A$) of a stretched flux tube should thin out ($A \propto 1/r$), failing to maintain a linear tension. We resolve this utilizing the \textbf{`Tensor Condensation'} mechanism derived in Part II \cite{YooPart2}. Bound by \textbf{Topological Flux Quantization}, the transverse geometric pressure locks the cross-sectional area $A$ as a constant non-dispersive soliton solution.

Integrating the covariant \textbf{`geometrical stress tensor'} ($S_{\mu\nu}$) yields a strictly linear potential energy $V(r)$. Thus, the QCD string tension ($\sigma$) \cite{Wilson1974} is directly derived from the geometric elasticity resisting topological unwinding:
\begin{equation}
V(r) = \int_{0}^{r} \mathcal{T}(x) \, dx = \left( \int_A S_{zz} \, dA \right) r \equiv \sigma r
\end{equation}

\subsection{Geomechanical Recovery of Asymptotic Freedom}
At extremely short distances ($r \to 0$), deep inside the core of the topological knot, the geometric stress tensor approaches the ideal fluid limit. As demonstrated in Part I via \textbf{`Topological Regularization,'} the divergent centrifugal phase kinetic energy exactly cancels the geometric potential well ($Q_s$). Consequently, the effective geometric tension drops asymptotically to zero ($\mathcal{T} \to 0$). This absolute cancellation allows the topological substructures to behave as asymptotically free wave-modes exclusively within the core, perfectly mirroring the empirical predictions of QCD \cite{Gross1973}.

\section{Emergent SU(3) Symmetry and Spinor Topology}

Mainstream physics attributes SU(3) symmetry to abstract color charges in internal 8-dimensional spaces. Recognizing the mathematical impossibility of directly mapping a 3D real space $SO(3)$ to a non-abelian complex $SU(3)$ Lie algebra, we propose an \textit{emergent} symmetry model.

\subsection{Emergence of SU(3) from 3D Topological Entanglement}
We postulate that the fundamental resonance axes ($x, y, z$) of the 3D spatial vibration serve as the physical basis for topological defects. The non-abelian SU(3) gauge field does not exist fundamentally but emerges macroscopically as the collective excitation of intertwined 3D topological knots. The geometric entanglement of three fundamental topological defects yields an effective resonance pattern that spans the state space:
\begin{equation}
\mathbf{3}_{\mathrm{topo}} \otimes \mathbf{3}_{\mathrm{topo}} \otimes \mathbf{3}_{\mathrm{topo}} \to \mathbf{10} \oplus \mathbf{8} \oplus \mathbf{8} \oplus \mathbf{1}
\end{equation}
Thus, Gell-Mann's Octet and Decuplet symmetries \cite{GellMann1964} are emergent, deterministic entanglement patterns arising from the collective geomechanics of the physical 3D tensor space.

\subsection{Spinor Topology and the Pauli Exclusion Principle}
As a spatial vortex rotates, it accumulates a topological Berry phase, imparting a strict \textbf{Spinor topology} to the knot. A complete topological unwinding requires a $4\pi$ ($720^\circ$) rotation of the phase plane. 

Classical wave destructive interference results merely in zero amplitude. However, in our geomechanical model, forcing two identical $4\pi$-twisted topological knots (fermions) into the exact same spatial coordinates demands an infinite topological deformation energy barrier to reconcile the orthogonal spinor phases. This absolute topological phase-exclusion is the deterministic, 3D physical origin of the \textbf{Pauli Exclusion Principle} \cite{Pauli1925}.

\section{Topological Jumps and Weak Interaction Decay}

\subsection{Strangeness as Topological Charge}
Acknowledging that 1-dimensional line integrals are insufficient to capture non-abelian topologies in simply connected spaces ($\pi_1(SU(3)) = 0$), we rigorously formalize ``Strangeness'' ($S$) as a non-trivial 3D Topological Charge ($Q_T$), evaluated via the Pontryagin index of the emergent gauge field tensor $F_{\mu\nu}$:
\begin{equation}
S \equiv Q_T = \frac{g^2}{32\pi^2} \int d^4x \, \mathrm{Tr}(F_{\mu\nu} \tilde{F}^{\mu\nu}) = n \quad (n \in \mathbb{Z})
\end{equation}
This mathematically proves that Strangeness must exist as discrete, quantized integers reflecting topological winding tension.

\subsection{Weak Decay via Quantum Instanton Tunneling}
Because $Q_T$ is a discrete integer ($n \in \mathbb{Z}$), it cannot decay continuously via classical viscous damping ($d\mathcal{W}/dt = -\gamma \mathcal{W}$). Instead, the unwinding of the topological knot (Weak Interaction) occurs via discontinuous \textbf{Topological Jumps}.

The spatial vacuum, possessing an intrinsic geometric chirality, selectively triggers these jumps in left-handed topologies. The decay rate ($\Gamma$) is mathematically governed by the quantum Instanton tunneling probability across the topological energy barrier ($S_{\mathrm{inst}}$):
\begin{equation}
\Gamma \propto \exp\left( -\frac{S_{\mathrm{inst}}}{\hbar} \right) \quad (\text{Parity-Violating Topological Jump})
\end{equation}
This resolves the mathematical contradiction, elucidating Parity Violation in weak interactions \cite{Wu1957} as discrete tunneling events governed by the chiral spatial fluid.

\subsection{Geomechanical Origin of W/Z Bosons}
When a highly tense topological flux tube abruptly tunnels to a lower winding state (unknotted), the sudden discrete release of torsional tension triggers a violent \textbf{`local tensor shockwave.'} Reusing the logic of fault rupture from Part III \cite{YooPart3}, the massive rest energy of the W and Z bosons corresponds exactly to the \textbf{`Threshold yield energy'} ($E_{\mathrm{crit}}$) released during this topological jump. The Higgs mechanism's mass generation \cite{Higgs1964} is thus elegantly replaced by the extreme \textit{Geometric Inertia} of the spatial medium resisting topological fracture.

\section{Conclusion}

In this concluding Part IV, we have systematically re-evaluated the epistemological necessity of abstract ``internal spaces'' in the Standard Model. By employing the principles of topological fluid dynamics, we proposed that hadrons are highly stable \textbf{Topological Tensor Knots}.

We overcame critical mathematical hurdles by abandoning continuous classical damping in favor of discontinuous \textbf{quantum instanton tunneling} to describe weak decays, ensuring the integrity of integer-valued topological charges. Furthermore, we resolved the $SO(3)-SU(3)$ dimensionality paradox by framing non-abelian gauge symmetries as \textit{emergent} collective excitations of 3D spatial tensor defects, preserving the rigorous mathematical structure of group theory within a 3D geometric ontology. 

The Pauli Exclusion Principle for fermions is deterministically derived from the infinite topological deformation energy barrier triggered by $4\pi$-twisted spinor phase planes. Color Confinement and Asymptotic Freedom are strictly derived from the conservation of fluid helicity (Hopf invariant) \cite{Moffatt1969} and the exact mathematical cancellation of nodal singularities (Part I) \cite{YooPart1, Gross1973}.

From the microscopic double-slit interference pattern (Part I) and the subatomic topological knots (Part IV), to the macroscopic condensation of dark matter (Part II) \cite{YooPart2} and the ultra-macroscopic filaments of the cosmic web (Part III) \cite{YooPart3}, nature does not change its laws across scales. Governed by a universal \textbf{`Fractal Scale Symmetry,'} the universe does not rely on abstract hidden dimensions, unphysical ad-hoc parameters, or dice-rolling. Particles, galaxies, and multiverses are merely navigating the deterministic curvatures and topological mechanics of a single, unified spatial vibration.

\vspace{1em}
\section*{ACKNOWLEDGMENTS}
\textbf{Note on Authorship:} This work was authored by Kwang Yong Yoo. The author acknowledges the use of AI-assisted tools (Gemini) for language refinement and simulation development support. This study is an independent research work conducted under the author's sole academic responsibility.

\begin{thebibliography}{99}
\bibitem{GellMann1964} M. Gell-Mann, A Schematic Model of Baryons and Mesons, \textit{Phys. Lett.} \textbf{8}, 214 (1964).
\bibitem{Wilson1974} K. G. Wilson, Confinement of Quarks, \textit{Phys. Rev. D} \textbf{10}, 2445 (1974).
\bibitem{Gross1973} D. J. Gross and F. Wilczek, Ultraviolet Behavior of Non-Abelian Gauge Theories, \textit{Phys. Rev. Lett.} \textbf{30}, 1343 (1973).
\bibitem{Moffatt1969} H. K. Moffatt, The Degree of Knottedness of Tangled Vortex Lines, \textit{J. Fluid Mech.} \textbf{35}, 117 (1969).
\bibitem{Pauli1925} W. Pauli, \"{U}ber den Zusammenhang des Abschlusses der Elektronengruppen im Atom mit der Komplexstruktur der Spektren, \textit{Z. Phys.} \textbf{31}, 765 (1925).
\bibitem{Wu1957} C. S. Wu \textit{et al.}, Experimental Test of Parity Conservation in Beta Decay, \textit{Phys. Rev.} \textbf{105}, 1413 (1957).
\bibitem{Higgs1964} P. W. Higgs, Broken Symmetries and the Masses of Gauge Bosons, \textit{Phys. Rev. Lett.} \textbf{13}, 508 (1964).
\bibitem{YooPart1} K. Y. Yoo, \textit{Mechanics of Spatial Vibration I: A Geomechanical Reinterpretation of Wave-Particle Duality and Phase Turbulence in Quantum Decoherence}, Zenodo. \url{https://doi.org/10.5281/zenodo.21206211} (2026).
\bibitem{YooPart2} K. Y. Yoo, \textit{Mechanics of Spatial Vibration II: Interaction of Macroscopic Gravity with Microscopic Vibration and a Unified Dark Fluid Model}, Zenodo. \url{https://doi.org/10.5281/zenodo.21233252} (2026).
\bibitem{YooPart3} K. Y. Yoo, \textit{Mechanics of Spatial Vibration III: Cosmological Spatial Plate Tectonics and Wave Interference Topology}, Zenodo. \url{https://doi.org/10.5281/zenodo.21258029} (2026).
\end{thebibliography}

\end{document}